\documentclass[10pt, a4paper]{article}
\usepackage{graphicx}
\graphicspath{{assets/}}
\usepackage{amsmath}

\title{My firstss LaTex document}
\author{Ghribi Ouassim Abdelmalek\thanks{Hello there}}
\date{April 2024}
\begin{document}
\maketitle
First document. This is a simple example, with no
extra parameters or packages included.
% This is a comment and should not appear

Some of the \textbf{greatest}
discoveries in \underline{science}
were made by \textbf{\textit{accident}}.

Some of the greatest \emph{discoveries} in science
were made by accident.

\textit{Some of the greatest \emph{discoveries}
    in science were made by accident.}

\textbf{Some of the greatest \emph{discoveries}
    in science were made by accident.}


\begin{figure}[h]
    \centering
    \includegraphics[width=0.75\textwidth]{alena-aenami-7p-m}
    \caption{Hello world}
    \label{fig:test-1}
\end{figure}

As you can see in figure \ref{fig:test-1} on page \pageref{fig:test-1}

\begin{itemize}
    \item First item
    \item Second item
\end{itemize}

\begin{enumerate}
    \item First enumerated item
    \item Second enumerated item
\end{enumerate}

In physics, the mass-energy equivalence is stated
by the equation $E=mc^2$, discovered in 1905 by Albert Einstein.

The mass-energy equivalence is described by the famous equation
\[ E=mc^2 \] discovered in 1905 by Albert Einstein.

In natural units ($c = 1$), the formula expresses the identity
\begin{equation}
    E=m
\end{equation}

Subscripts in math mode are written as $a_b$ and superscripts are written as $a^b$. These can be combined and nested to write expressions such as

\[ T^{i_1 i_2 \dots i_p}_{j_1 j_2 \dots j_q} = T(x^{i_1},\dots,x^{i_p},e_{j_1},\dots,e_{j_q}) \]

We write integrals using $\int$ and fractions using $\frac{a}{b}$. Limits are placed on integrals using superscripts and subscripts:

\[ \int_0^1 \frac{dx}{e^x} =  \frac{e-1}{e} \]

Lower case Greek letters are written as $\omega$ $\delta$ etc. while upper case Greek letters are written as $\Omega$ $\Delta$.

Mathematical operators are prefixed with a backslash as $\sin(\beta)$, $\cos(\alpha)$, $\log(x)$ etc.

\section{First example}

The well-known Pythagorean theorem \(x^2 + y^2 = z^2\) was proved to be invalid for other exponents, meaning the next equation has no integer solutions for \(n>2\):

\[ x^n + y^n = z^n \]

\section{Second example}

This is a simple math expression \(\sqrt{x^2+1}\) inside text.
And this is also the same:
\begin{math}
    \sqrt{x^2+1}
\end{math}
but by using another command.

This is a simple math expression without numbering
\[\sqrt{x^2+1}\]
separated from text.

This is also the same:
\begin{displaymath}
    \sqrt{x^2+1}
\end{displaymath}

\ldots and this:
\begin{equation*}
    \sqrt{x^2+1}
\end{equation*}

\end{document}